\documentclass[11pt]{article}
\usepackage{amsmath,amssymb,amsthm,booktabs,hyperref}
\newtheorem{theorem}{Theorem}
\newtheorem{lemma}[theorem]{Lemma}
\newtheorem{proposition}[theorem]{Proposition}
\title{The intrinsic centre and finite gluing interfaces}
\author{Human Compression project}
\date{}
\begin{document}
\maketitle

\paragraph{Scope.}
This concerns nonbipartite quintic graphs of order $32$, girth four and
diameter eight. It is a written extension of the independently reviewed
twenty-four-row classification~\cite{rows}. It gives an intrinsic interface for ordinary
identification and a larger right-half domain containing the earlier
singleton domain. These new statements await independent review.
No graph, Source, matrix, TeX or Lean execution was performed in preparing
this note, and no numerical census is asserted.

We retain the notation of~\cite{rows}: the midpoint is $p_4$, the two
sides are $S,T$, the ports are $B_S,B_T,B_0$, the neutral isolated middle
vertices number $z$, and $j=|B_0|$. Every constructed graph has diameter
eight and every vertex is within four of $p_4$. Its centre therefore
consists precisely of the vertices of eccentricity four. Write
$X=\{p_4\}\sq M_{\rm iso}$; these $z+1$ vertices are false twins, with
common neighborhood $B$.

\section{Determining the centre}

\begin{lemma}\label{restrict}
Every central vertex lies in $X$, in $B_0$, or among the neutral middle
leaves in a middle-edge row. No other vertex is central.
\end{lemma}
\begin{proof}
Every endpoint row, side vertex or side port is within three of an
endpoint. The same holds for a middle vertex with a side neighbor.
If $d(v,p_0)\leq3$, then $d(v,p_8)\geq5$; the opposite statement is
identical. The twenty-four-row classification leaves only the named
possibilities outside those types. All members of $X$ are central
because they are false twins of $p_4$.
\end{proof}

\begin{proposition}\label{nomiddle}
In a row with no middle--middle edge, the centre is $X$ together with
all of $B_0$ if both sides have a middle neighbor. In the sole row
where $B_0$ is nonempty but a side has no middle neighbor, namely the
internal-edge row $j=1$, no member of $B_0$ is central. Its centre is
the four-vertex false-twin class $X$.
\end{proposition}
\begin{proof}
All vertices of $B_0$ are false twins with neighborhood
$\{p_4\}\sq M$, so they have the same eccentricity. Suppose first that
both sides have middle neighbors and let $y\in B_0$. Every middle vertex
and $p_4$ is adjacent to $y$. Every side port has a middle neighbor and
is therefore within two of $y$. This follows immediately from $z>0$;
the only row with $z=0$ is $Q_7$, where the displayed middle--port table
gives the same assertion.

On a side without an internal edge, a vertex with a middle neighbor is
within two of $y$, and a vertex with a port neighbor is within three.
Let $Z$ be the union of these vertices. Its column deficits sum to five
and are at most four, so $|Z|\geq2$. Each endpoint four-subset meets
$Z$ and is within four of $y$. A side vertex outside $Z$ meets all five
endpoint rows, including a row incident with a middle-neighbor vertex;
it is within four of $y$ as well. The endpoint itself has distance at
most four via a middle neighbor and its side neighbor.

On a side with an internal edge, this is one of the $j=2,3$ rows.
The middle neighbors all meet one endpoint of that edge, giving it
distance two from $y$. The side ports meet its other endpoint, giving
that endpoint distance at most three. Every endpoint row contains one
of these endpoints and all three side isolates; all these vertices
again have distance at most four from $y$. The other side has no
internal edge and was handled above. The remaining vertices of $B_0$
are at distance two from $y$. Thus $y$ has eccentricity at most four
and is central.

In the internal-edge row $j=1$, the side $S$ has no middle neighbor,
$z=3$, and the one nonisolated middle vertex $u$ lies on $T$.
The neighborhood of $y\in B_0$ is $X\sq\{u\}$. Each member of $X$
has distance four from the $S$ endpoint. The vertex $u$ is within
three of the $T$ endpoint, and hence has distance at least five from
the $S$ endpoint. Every path starting at $y$ therefore has length at
least five before reaching the $S$ endpoint, and a path through $X$
has length exactly five. Thus $y$ is not central. Lemma~\ref{restrict}
completes the proof.
\end{proof}

\begin{proposition}\label{middle}
In a middle-edge row, all $r$ neutral middle leaves are central. They
form a second false-twin class $V$, disjoint from $X$. The full centre
is the independent set $X\sq V$.
\end{proposition}
\begin{proof}
Retain the degree-four middle centre $u\in M_S$ and its sole port
neighbor $a\in B_S$. Each neutral leaf $v$ has neighborhood
$(B\setminus\{a\})\sq\{u\}$, so these leaves are false twins.
It suffices to show one such $v$ has eccentricity four.

The vertex $u$ has distance one from $v$ and $a$ has distance two.
All other ports have distance one. The vertices of $X$, all other
neutral leaves, and the additional middle vertex on $T$ when present
have distance at most two. Every star leaf on $S$ has distance two.
A side vertex adjacent to $a$ has distance at most three; a vertex
adjacent to the other $B_S$ port, when present, has distance two.
A remaining side vertex has no residual edge or port incidence and
meets all five endpoint rows. Choosing a row containing a star leaf
puts this vertex within four of $v$.

If $r\leq2$, every endpoint four-subset on the five-vertex side $S$
meets a star leaf and is within three of $v$. If $r=3$, a row missing
the sole star leaf contains all four other side vertices, hence a
neighbor of $a$; it is within four of $v$. The endpoint is within
four through a row containing a star leaf.

All ports on $T$ are adjacent to $v$. The retained six-vertex radius
lemma applies with $v$ as hub in the $q=3$ rows. In the $q=4$ rows,
the five-vertex radius lemma applies: the possible additional middle
vertex on $T$ is within two of $v$. Thus every vertex on that side
and its endpoint rows is within four, and its endpoint is within four
through a side-port neighbor. This proves the required eccentricity.
Lemma~\ref{restrict} excludes all other vertices. Neither $X$ nor $V$
has internal edges and there are no edges between them. Their neighborhoods
differ by $a$ versus $u$, so they are distinct full-graph twin classes.
\end{proof}

These propositions give the following intrinsic centre types. Here $E_w$
denotes one independent false-twin class of size $w$, and $E_{a,b}$
denotes two distinct independent false-twin classes of sizes $a,b$ with
no edges between them.
\[
\begin{array}{c|c|c}
\text{row family}&\text{centre}&\text{component orders after deletion}\\\hline
q=2\text{, no middle edges}&E_2&(14,16)\\
q=3\text{, no middle edges, degree }d&E_3&(17-d,12+d)\\
q=3\text{, middle edges}&E_{3-r,r}&(13,16)\\
q=4\text{, middle edges, opposite side edge}&E_{4-r,r}&(14,14)\\
q=4\text{, middle edges, opposite middle edge}&E_{3-r,r}&(13,16)\\
Q_1,Q_2&E_3&(13,16),(14,15)\\
Q_3&K_{3,2}&(13,14)\\
Q_4&K_{2,1}&(13,16)\\
Q_5&K_{2,2}&(14,14)\\
Q_6&K_{2,3}&(13,14)\\
Q_7&K_{1,2}&(13,16)\\
\text{internal side edge }j=1&E_4&(14,14)\\
\text{internal side edge }j=2,3&K_{5-j,j}&(14,13)
\end{array}
\]
In each complete bipartite centre, its two parts are its full-graph
false-twin classes. The table does not assert that different rows give
different ordinary graphs.

\begin{theorem}\label{decompose}
Deleting the intrinsic centre leaves exactly two connected components,
each of order at most sixteen. The centre has at most five vertices and
at most two full-graph false-twin classes.
\end{theorem}
\begin{proof}
The centre descriptions were proved above. Each side and its endpoint
rows lie in a connected component: the endpoint meets all five rows,
and every side vertex has positive incidence with them. Side ports
also meet that side. Every remaining middle vertex has a side neighbor.
In middle-edge rows the centre leaves are removed, while the vertex
$u$ and its port $a$ stay on the same side. The only noncentral member
of $B_0$ is in the internal-edge $j=1$ row; after deleting $X$ it has
the sole neighbor $u\in M_T$, and therefore stays on $T$.
There is no edge between the two resulting regions, and neither is
empty. Hence they are exactly the two components.

Their orders follow by counting the five side vertices, five endpoint
rows, one endpoint, their side ports and noncentral middle vertices;
the six-vertex rows add one extra side vertex. For example a $q=3$
no-middle-edge row has orders $5+5+1+1+(5-d)=17-d$ and
$6+5+1+d=12+d$. In a complete bipartite centre row the orders are
$11+m_S+b_S$ and $11+m_T+b_T$. The other entries follow from the
explicit port tables. These are all the entries displayed, and all
are at most sixteen. The centre class sizes also follow from the
proved descriptions.
\end{proof}

\section{Ordinary gluing without a choice of geodesic}

Let $C$ be the intrinsic centre. Give each full-graph false-twin class
of $C$ one coordinate. An outside vertex is adjacent to either all or
none of that class; record these choices as a one- or two-bit boundary
color. Record also each centre class size and whether the two classes
are joined. Within a false-twin class there are no edges.

For each of the two components of $G-C$, take any complete canonical
code for its ordinary graph with these vertex colors. Such a code is
available by a finite search over color-preserving vertex orders; a
bounded implementation can use an exact refinement search instead.
Sort the two component codes. If the centre classes have distinct
sizes, order them by size. If they have the same size, take the lesser
result from the two possible coordinate orders. There is no further
alignment of individual centre vertices: their permutations act trivially
on every adjacency to the components.

\begin{theorem}\label{canonical}
The weighted centre data and the sorted pair of component codes,
minimized over the indicated at most two class orders, form a complete
invariant of the ordinary uncolored graph in this class.
\end{theorem}
\begin{proof}
An ordinary isomorphism preserves eccentricity, hence $C$, and preserves
the full-neighborhood equivalence relation on it. It preserves class
sizes and the centre adjacency bit. It maps the two components and
their boundary colors according to one of the permitted centre-class
orders. Thus it preserves the stated code.

Conversely the code reconstructs the complete graph. Expand each centre
class into its specified independent set, join the classes if indicated,
take the two coded colored components, and join each outside vertex to
the centre classes selected by its color. Equal codes therefore give a
common reconstructed graph and explicit full-vertex isomorphisms to it.
This proves completeness. A permutation among equal-role centre vertices
or an exchange of equal-order components introduces no unrecorded choice.
\end{proof}

This permits canonicalization of a side leg before crossing side
catalogues. Every local component isomorphism preserving its boundary
colors extends by the identity on the labelled centre. Two independent
leg isomorphisms therefore glue to an isomorphism of the full graph.
Marked side repetitions may be removed at this stage. If the two centre
classes have equal size, each leg stores both color-coordinate versions;
the global choice must be simultaneous for the two legs. Independent
minimization under separate swaps would lose their correlation.

\section{The complete two-centre class}

The earlier right-half domain~\cite{right} has an endpoint $p$, five
endpoint rows $A$, seven side vertices $D$ with three internal edges,
and three intrinsic degree-three ports $P$. The four possibilities for
$D$ are $P_4\sq3K_1$, $K_{1,3}\sq3K_1$,
$P_3\sq K_2\sq2K_1$, and $3K_2\sq K_1$.
Endpoint and port rows are independent four- and three-subsets, with
\[
 \alpha(v)+\beta(v)+d_D(v)=5.
\]
The earlier singleton domain required $\alpha(v)>0$.

\begin{lemma}\label{zero}
If that condition is relaxed to $\alpha(v)\geq0$, there is at most one
zero column. If it exists, its degree in $D$ is two or three and it is
exactly the exceptional nonisolated middle vertex in the two $q=2$ rows.
Every such right half is connected.
\end{lemma}
\begin{proof}
A zero column requires $d_D(v)+\beta(v)=5$ with $\beta(v)\leq3$;
therefore $d_D(v)\geq2$. The three-leaf star and $P_3\sq K_2$ each
have only one such possible vertex, and the matching has none. The only
two candidates in $P_4$ are adjacent. If both had zero columns, both
would belong to all three port subsets, contradicting their independence.
Thus there is at most one.

Removing that vertex leaves six positive side columns. If its degree
is three, the remaining side graph is empty and two ports meet it;
the other port does not. If its degree is two, all three ports meet
it, and there is one remaining internal side edge. These are precisely
the $q=2$ side incidences in~\cite{rows}; the opposite six-vertex side
is the retained $P_3$ leg with two ports of arity three. Every positive
side column meets an endpoint row and is connected to $p$. Each port
has three side neighbors and at most one of them can be a zero column,
so each is connected to $p$. The zero column, if present, meets at least
two ports, and is connected as well.
\end{proof}

\begin{proposition}\label{strictdisjoint}
A right half with a zero-column presentation admits no positive-column
presentation at any other endpoint. The positive and zero-column domains
are therefore disjoint as ordinary unrooted half classes.
\end{proposition}
\begin{proof}
Use the $q=2$ notation: the degree-five vertices are $p$, the five rows
$A$, six ordinary side vertices $S$, and the zero-column vertex $u$.
The ports are intrinsically degree three. A possible new endpoint must
avoid the ports, have no port-neighbor in its neighborhood, and give
positive incidence to every remaining side vertex~\cite{right}.

The vertex $p$ fails because of $u$; $u$ itself meets at least two ports.
Consider a side vertex $v$. If it meets a port, it is not an endpoint.
If it meets $u$, its neighborhood contains the port-neighbor $u$ and
again fails. If it meets neither, its only possible residual neighbor
is along the one internal side edge. If $vw$ is that edge and both
vertices had zero port incidence, their endpoint-row incidence sum
would be at least seven, impossible across five independent rows.
Thus $w$ meets a port, again excluding $v$. With no such residual edge,
$v$ has precisely the neighborhood $A$ of $p$ and still leaves $u$
with zero endpoint incidence.

Finally consider a row vertex $a\in A$. Its new endpoint neighborhood
is $\{p\}$ together with four vertices of $S$. The two omitted side
vertices can meet this neighborhood only through the at most one
internal side edge. At most one of them can do so. Hence the new
presentation again has a zero side column. This exhausts all thirteen
degree-five vertices and proves the claim.
\end{proof}

For the enlarged half domain, the ordinary code is still obtained by
all eligible endpoints and all side isomorphisms, now dropping only the
strict positive-column test. Every candidate still gives exactly three
internal side edges by the same edge count, and all four side types
remain unchanged. The bound of thirteen endpoint choices and forty-eight
side maps, hence at most $624$ candidate tuples, is unchanged. Its
ordinary completeness proof is the same full reconstruction argument.
The existing strict kernel and its public reproduction are not modified
by this proposal; an extension needs its own reviewed Source identity.

\begin{theorem}\label{twocentre}
The complete nonbipartite class in this setup whose centre is a two-vertex false-twin
class consists of the singleton stratum together with the two $q=2$
nonisolated rows. Its ordinary count is six times the ordinary unrooted
count of the enlarged nonnegative-column right-half domain. This is a
structural identity, not a numerical evaluation.
\end{theorem}
\begin{proof}
The centre table excludes every other nonisolated row. The retained
all-isolated middle cases~\cite{singleton} have centre size three or four; the singleton
case has centre size two. The empty middle case is impossible.
Deleting the centre gives connected components of orders fourteen and
sixteen. Their side-port sets are intrinsic degree-three vertices,
numbering two and three. The fourteen-vertex side is one of the six
retained unrooted forms in both the singleton and $q=2$ cases.
The sixteen-vertex side is exactly the enlarged domain by Lemma~\ref{zero}.
Conversely every such pair glues to a target graph by the respective
retained singleton or reviewed $q=2$ construction and its radius proof.

The two component orders distinguish them intrinsically, and each
gluing joins every port to both centre vertices. Thus there are no
alignment choices. An ordinary isomorphism preserves the centre and
the two components, and conversely component isomorphisms glue. The
claimed product follows. Proposition~\ref{strictdisjoint} records the
additional disjointness from the earlier strict half domain.
\end{proof}

\paragraph{Remaining work.}
All new centre and enlarged-domain claims above require independent
review before activation. Exact colored-component canonicalization,
bounded side-catalogue enumeration, complete side products and ordinary
union across the retained strata still require reviewed Source and actual
Root execution. No expected count may be used as a target or stopping
rule, and this note supplies no whole manuscript or new conjecture credit.

\begin{thebibliography}{9}
\bibitem{rows} Human Compression project, \emph{Nonisolated middle vertices
in quintic graphs of diameter eight}, frozen source SHA256
\texttt{522af65bc2bb25d0d09ad68b1c047f2806832230e8452d388ead3483d84f6384},
with independent review by the separate agent.
\bibitem{right} Human Compression project, \emph{Finite enumeration of the
right halves in the singleton middle case}, frozen source SHA256
\texttt{1588351c747e7e44c80d1b5736c5ac6addfd2886a2ab283301ce1e848b16229a}.
\bibitem{singleton} Human Compression project,
\emph{The empty and singleton middle cases for quintic graphs}, frozen
source SHA256
\texttt{76a5c09c2e3f52536c022572bd95ba611392ad38527c0f2ee40c1e886e9d9eb2},
with independent review SHA256
\texttt{026a34a5c0d80143005551a931a6eb5d83c899f29a048ed52aa72e075c5b6e63}.
\end{thebibliography}
\end{document}
